% ============================================================
% 07-conclusion.tex — V13 Conclusion (condensed, filled)
% ============================================================
\section{Conclusion and Future Work}
\label{sec:conclusion}

This study developed and evaluated \system{}, a supervisor-routed multi-agent RAG framework for heterogeneous university counseling, through architecture-level comparison, mechanistic ablation, and tool-ambiguity stress testing.

Three principal findings address the research questions. For RQ1, the architecture-level comparison indicates that supervisor-routed multi-agent decomposition achieves higher End-to-End Success (81.0\% vs.\ 73.0\%) and 41.4\% fewer input tokens than a functionally matched single-agent baseline, with an associated latency trade-off. For RQ2, the ablation reveals that specialist policies ($-10.0$\,pp E2E) and heterogeneous evidence allocation ($-12.3$\,pp Source Recall) are the most impactful components, with structured-financial and cross-domain queries exhibiting the largest degradation when graph grounding is removed. For RQ3, the Routed Specialist preserves 88.0\% E2E pass at 51 tools while the Single Agent degrades to 83.3\% and the Global Top-10 drops to 78.7\%.

These findings are consistent with the interpretation that supervisor-routed bounded specialization---comprising domain-specific policies, scoped tool access, representation-specific evidence paths, and deterministic computation---can support heterogeneous institutional counseling under the evaluated conditions, with measurable trade-offs in coordination latency and token cost.

Future work should (1)~reduce routing latency through parallel candidate dispatch and prompt compression; (2)~extend evaluation to multi-institutional and bilingual benchmarks with naturally evolving tool registries; (3)~conduct human evaluation with multiple raters to complement automated metrics; and (4)~perform field trials assessing student trust and advisory efficacy in live counseling settings.
